\section{Обучение и оценка качества}

\subsection{Обучающая разметка}

Для каждой позиции обучающей выборки сначала выполняются все \(K\) внутренних шагов. Затем
вычисляются потери \(L_{t,0},\ldots,L_{t,K}\) и оракульное время \(\tau_t^\star\) по
формуле~\eqref{eq:oracle_stopping_time}. После этого обучается модель распределения \(\widehat
p_{t,k}\).

Если используется точечная разметка, целевая функция имеет вид
\begin{equation}
  \mathcal L_{\mathrm{point}}
  =-\sum_t\sum_{k=0}^{K}\log \widehat p_{t,k}(\tau_t^\star).
  \label{eq:point_loss}
\end{equation}
Если используется сглаженная разметка, применяется функция~\eqref{eq:distribution_loss}. Оба
варианта нужно сравнить, поскольку точечная разметка проще, а сглаженная лучше отражает близкие по
риску шаги.

\subsection{Качество языкового прогноза}

Основная метрика качества прогноза --- средняя отрицательная логарифмическая потеря при
фактическом времени остановки:
\begin{equation}
  \overline L=\frac{1}{N}\sum_{t=1}^{N}-\log Q_{t,\tau_t}(X_{t+1}).
  \label{eq:mean_nll}
\end{equation}
Дополнительно фиксируется доля совпадений верхнего варианта с правильной языковой единицей:
\begin{equation}
  \mathrm{Acc}=\frac{1}{N}\sum_{t=1}^{N}
  \ind\left\{\argmax_{a\in\A}Q_{t,\tau_t}(a)=X_{t+1}\right\}.
  \label{eq:accuracy}
\end{equation}
Эти метрики нужны вместе: логарифмическая потеря чувствительна к качеству распределения, а доля
совпадений показывает поведение верхнего варианта.

\subsection{Качество предсказания времени}

Для самой модели времени используются два показателя. Первый --- отрицательная логарифмическая
потеря по \(\tau^\star\):
\begin{equation}
  \mathrm{NLL}_{\tau}=-\frac{1}{N(K+1)}\sum_t\sum_{k=0}^{K}\log\widehat p_{t,k}(\tau_t^\star).
  \label{eq:tau_nll}
\end{equation}
Второй --- средняя абсолютная ошибка математического ожидания:
\begin{equation}
  \mathrm{MAE}_{\tau}=\frac{1}{N(K+1)}\sum_t\sum_{k=0}^{K}
  \left|\sum_{j=0}^{K}j\widehat p_{t,k}(j)-\tau_t^\star\right|.
  \label{eq:tau_mae}
\end{equation}

\subsection{Вычислительная стоимость}

Стоимость измеряется средним числом внутренних шагов
\begin{equation}
  \overline\tau=\frac{1}{N}\sum_{t=1}^{N}\tau_t
  \label{eq:mean_tau}
\end{equation}
и долей экономии относительно полного запуска:
\begin{equation}
  \mathrm{Save}=1-\frac{\overline\tau}{K}.
  \label{eq:saving}
\end{equation}
Если учитывается начальный шаг \(0\), в знаменателе можно использовать \(K+1\); в экспериментах
нужно зафиксировать один вариант и не менять его между таблицами.

\subsection{Словарь обозначений метрик}\label{sec:metrics_legend}

Ниже собраны основные величины, встречающиеся в тексте и таблицах; формулы для \(\overline L\),
\(\mathrm{Acc}\), \(\mathrm{NLL}_\tau\), \(\mathrm{MAE}_\tau\), \(\overline\tau\), \(\mathrm{Save}\)
даются выше в разделе.

\textbf{Качество языка и времени остановки.}
\(\overline L\) --- средняя отрицательная логарифмическая потеря~\eqref{eq:mean_nll} при
выбранном реализуемом \(\tau_t\); \(\mathrm{Acc}\) --- доля позиций, на которых аргмаксим
распределения совпал с истинным следующим символом~\eqref{eq:accuracy}; \(\mathrm{NLL}_\tau\) ---
средняя перекрёстно-энтропийная потеря для меток \(\tau_t^\star\)~\eqref{eq:tau_nll};
\(\mathrm{MAE}_\tau\) --- средняя абсолютная ошибка оценки \(\E[\tau^\star\mid\Hh_{t,k}]\)
относительно \(\tau_t^\star\)~\eqref{eq:tau_mae}.

\textbf{Стоимость вычислений.}
\(\overline\tau\) и \(\mathrm{Save}\) --- среднее число внутренних шагов и относительная
экономия~\eqref{eq:mean_tau}, \eqref{eq:saving} относительно полного прогона на \(K\) шагов.

\textbf{Таблицы результатов (WikiText и ARC).}
\(\mathrm{tok}_{\mathrm{best}}\) --- достигнутая токеновая точность (топ-1) в точке остановки при
выбранной политике; \(\mathrm{top5}_{\mathrm{best}}\) --- доля попаданий истинного токена в пять
лучших кандидатов (только для языковой постановки). В ARC дополнительно фиксируется
\(\mathrm{NLL}_\tau\) в точке остановки как вспомогательная величина качества прогноза.

\textbf{Сглаживание графиков (EMA).}
В подписях к рисункам обучения аббревиатура EMA означает \emph{exponential moving average}: то же
самое, что экспоненциальное сглаживание отображаемой величины по итерациям; это способ
визуализации траектории, а не отдельная метрика математической модели.